\documentclass[UTF8,a4paper,12pt]{ctexart}

\usepackage{geometry}
\usepackage{setspace}
\usepackage{graphicx}
\usepackage{booktabs}
\usepackage{array}
\usepackage{longtable}
\usepackage{xcolor}
\usepackage{hyperref}

\geometry{left=2.8cm,right=2.8cm,top=2.6cm,bottom=2.6cm}
\setstretch{1.35}
\hypersetup{
    colorlinks=true,
    linkcolor=black,
    urlcolor=black,
    citecolor=black
}

\newcommand{\swufelogo}{\includegraphics[width=10cm]{swufe_logo.jpg}}

\title{QuantMind Agent：基于多 Agent 的股票辅助分析系统设计与实现}
\author{孙溥~42311043、夏羽欣~42311160、苏比~42311185\\西南财经大学 计算机与人工智能学院}
\date{2025--2026学年第二学期}

\begin{document}
\begin{titlepage}
\centering
\vspace*{0.9cm}
\noindent\makebox[\textwidth][c]{\raisebox{-9cm}[0pt][0pt]{\swufelogo}}\par
\vspace{-\baselineskip}
{\zihao{0}\heiti 西南财经大学\par}
\vspace{0.4cm}
{\zihao{1} Southwestern University of Finance and Economics\par}
\vspace{1.0cm}
{\zihao{0}\heiti 课程论文\par}
\vspace{3.5cm}
\renewcommand{\arraystretch}{1.8}
\begin{tabular}{>{\heiti}r p{9cm}}
学年学期： & 2025--2026--2 \\
课程名称： & 算法交易 \\
论文题目： & QuantMind Agent：基于多 Agent 的股票辅助分析系统设计与实现 \\
小组成员： & 孙溥~42311043、夏羽欣~42311160、苏比~42311185 \\
学院： & 计算机与人工智能学院 \\
年级专业： & 23级计算机科学与技术 \\
\end{tabular}
\vfill
{\zihao{4} \today\par}
\end{titlepage}
\renewcommand{\arraystretch}{1.15}

\tableofcontents
\newpage

\begin{abstract}
随着金融市场信息来源不断增多，投资者在进行股票分析时需要同时关注行情走势、技术指标、新闻信息、基本面质量、市场舆情和风险水平。传统单一规则模型虽然实现简单，但难以综合处理多维信息；而直接使用单个大语言模型进行判断，也可能存在信息混杂、推理过程不清晰和结果不稳定等问题。为此，本文设计并实现了 QuantMind Agent，一个面向股票辅助分析的多 Agent 系统。系统参考 TradingAgents 多 Agent 金融交易框架的思想，但结合课程项目规模和实际部署需求，采用可解释的串行协作流程，设置技术分析 Agent、新闻分析 Agent、基本面分析 Agent、舆情分析 Agent、多头研究员 Agent、空头研究员 Agent、研究经理 Agent、风险控制 Agent 和交易决策 Agent 九个角色。系统后端基于 Python 与 FastAPI 实现，前端使用 HTML、CSS 和 JavaScript 构建交互页面，支持用户输入 A 股或美股名称及代码，并通过服务端事件推送展示分析进度。系统能够获取行情数据、新闻数据和基本面数据，计算 MA5、MA10、最新价格、成交量变化等指标，并结合新闻情绪、财务指标、舆情热度、多空研究结论和风险控制规则生成 BUY、HOLD、WAIT、SELL 等辅助决策结果。项目已部署至 Zeabur 平台，并在 GitHub 开源。实验结果表明，系统能够完成股票输入识别、Agent 流程执行、分析结果展示和异常数据安全回退等功能，具有一定的可解释性和实用价值。需要说明的是，本文系统仅用于课程研究和学习展示，不构成任何真实投资建议。

\noindent\textbf{关键词：}多 Agent；大语言模型；股票分析；辅助交易决策；风险控制；数据可视化
\end{abstract}

\newpage

\section{引言}

\subsection{研究背景}

股票市场受到公司基本面、市场情绪、宏观政策、行业新闻和资金行为等多种因素影响。普通投资者在进行股票分析时，通常需要同时查看价格走势、成交量变化、均线指标、新闻事件和风险水平。由于信息来源分散、数据类型多样，仅依靠人工阅读和经验判断容易出现分析效率低、判断依据不完整等问题。

近年来，大语言模型在文本理解、摘要生成和推理分析方面表现出较强能力，为金融辅助分析提供了新的思路。与传统量化模型相比，大语言模型可以更自然地处理新闻、公告、研报等非结构化文本；与黑箱深度学习模型相比，基于语言模型的 Agent 系统能够输出较为清晰的自然语言解释。TradingAgents 等研究进一步提出，可以通过多个具有不同角色的 Agent 模拟真实交易机构中的分析师、交易员和风险控制团队，从而增强金融决策过程的结构化和可解释性。

本文基于上述背景，设计并实现了 QuantMind Agent 系统。该系统不是自动化实盘交易平台，而是一个面向课程研究和学习展示的股票辅助分析系统。系统在线演示地址见\href{https://quantmind-agent.zeabur.app/}{QuantMind Agent 在线演示平台}，项目源代码地址见\href{https://github.com/sunpu24/QuantMind-Agent}{GitHub 项目仓库}。

\subsection{研究意义}

本文研究具有以下意义。第一，从应用角度看，系统将行情数据、新闻信息、技术指标、基本面指标、市场舆情和风险控制整合到统一流程中，可以帮助用户更直观地理解一只股票当前的技术面、消息面、基本面和情绪面情况。第二，从系统设计角度看，多 Agent 架构将复杂股票分析任务拆分为若干相对独立的模块，每个 Agent 负责一个明确任务，有利于提升系统的可维护性和可扩展性。第三，从决策解释角度看，系统新增多头研究员、空头研究员和研究经理 Agent，使最终决策不仅依赖单一分析报告，还能够呈现多空观点对比和证据汇总过程。第四，从数据可视化角度看，系统通过网页输入、进度条和结果卡片展示分析过程，使抽象的 Agent 决策流程变得更加直观。第五，从风险控制角度看，系统在数据不可用或回退到 Mock 数据时，会自动将最终交易建议降级为 WAIT，并将仓位设置为 0，从而避免基于不可靠数据给出过度明确的买卖建议。

\subsection{研究内容}

本文围绕 QuantMind Agent 系统展开，主要研究内容包括以下几个方面：

第一，设计股票辅助分析系统的整体架构。系统包括前端交互页面、FastAPI 后端服务、股票输入识别模块、行情数据模块、新闻数据模块、多 Agent 工作流和结果展示模块。

第二，实现多层次多 Agent 分析流程。系统设置技术分析 Agent、新闻分析 Agent、基本面分析 Agent、舆情分析 Agent、多头研究员 Agent、空头研究员 Agent、研究经理 Agent、风险控制 Agent 和交易决策 Agent 九个核心角色，按照固定串行流程依次执行。

第三，实现行情、新闻与基本面数据处理。系统支持 AkShare、Tushare、Alpha Vantage 等数据源，同时保留 Mock 或空数据回退方案，用于测试和异常保护。

第四，实现风险控制与安全提示机制。系统通过风险评分、建议仓位、止损比例和 Mock 数据保护规则，降低辅助决策结果的误导风险。

第五，实现 Web 可视化展示。系统通过网页展示股票输入、Agent 执行进度、最终交易决策和各 Agent 的分析详情。

\subsection{论文结构安排}

本文共分为五个部分。第一部分为引言，介绍研究背景、研究意义、研究内容和论文结构。第二部分为文献综述，介绍大语言模型在金融领域的应用、多 Agent 协作框架以及技术分析和新闻情绪分析相关研究。第三部分为研究设计，说明 QuantMind Agent 的系统需求、总体架构、Agent 流程、数据处理和 Web 展示设计。第四部分为实验分析，介绍实验环境、实验流程、功能测试和案例分析。第五部分为结论与展望，总结本文工作并说明系统不足和后续改进方向。

\section{文献综述}

\subsection{大语言模型在金融领域的应用}

大语言模型能够处理大量自然语言文本，并在问答、摘要、分类和推理任务中表现出较强能力。在金融领域，相关研究主要集中在金融问答、财报分析、新闻摘要、情绪识别和金融信息抽取等方面。例如 BloombergGPT、FinGPT、PIXIU 等模型或框架通过金融语料训练、指令微调或金融任务评测，提升模型对金融术语、财务指标和市场文本的理解能力。

这类研究表明，大语言模型可以作为金融分析助手，为用户提供信息整理和辅助判断能力。不过，仅将大语言模型作为单一问答工具仍存在一定局限。例如，模型可能缺少实时行情数据，或者在面对多源信息时难以保持稳定的结构化推理。因此，将大语言模型与外部数据源、规则系统和 Agent 流程结合，是提升金融应用可靠性的一种重要方向。

\subsection{基于大语言模型的股票分析研究}

在股票分析任务中，大语言模型不仅可以用于阅读新闻和公告，也可以结合行情指标、财务信息和市场情绪生成辅助判断。已有研究通常可以分为新闻驱动、推理驱动和回测反馈驱动等类型。新闻驱动方法主要使用新闻标题、正文或社交媒体信息判断市场情绪；推理驱动方法强调模型对不同信息的综合解释；回测反馈驱动方法则尝试根据历史交易结果优化模型输出。

需要注意的是，股票市场具有较高不确定性，任何模型都难以保证预测准确。因此，在实际系统设计中，应避免将模型输出直接包装成确定性的投资建议，而应将其定位为辅助分析结果。本文系统在最终报告中明确加入免责声明，并通过 WAIT、仓位为 0、风险提示等方式降低误导风险。

\subsection{多 Agent 协作框架研究}

多 Agent 系统通过为不同 Agent 分配不同角色，将复杂任务拆解为多个子任务。TradingAgents 提出了一个面向金融交易的多 Agent 框架，模拟真实交易机构中的分析师团队、研究员团队、交易员、风险管理团队和基金经理审批流程。该框架的核心思想是让不同 Agent 从不同角度收集信息、展开讨论、控制风险并生成交易决策。

本文系统参考了 TradingAgents 的角色分工思想，但没有完全复制其复杂结构。考虑到课程项目的实现难度、部署成本和展示需求，QuantMind Agent 采用轻量化的固定串行流程，将系统划分为技术分析、新闻分析、基本面分析、舆情分析、多头研究、空头研究、研究经理、风险控制和交易决策九个 Agent。与旧版只包含基础分析和决策模块的流程相比，当前系统进一步引入轻量级多空研究与研究经理汇总机制，能够在保持实现复杂度可控的前提下展示多 Agent 分工协作和观点整合过程。

\subsection{技术分析与新闻情绪分析研究}

技术分析是股票分析中的常见方法，主要根据价格、成交量和技术指标判断市场趋势。常见指标包括移动平均线、MACD、RSI、KDJ、布林带等。本文系统当前实现了较为基础但易解释的技术指标，包括 MA5、MA10、最新价格和成交量变化率。系统根据最新价格与均线之间的关系判断技术面偏强、偏弱或中性。

新闻情绪分析则关注新闻事件对市场预期和投资者情绪的影响。公司公告、行业政策、宏观经济变化和负面事件都可能影响股票价格。大语言模型适合处理新闻标题和摘要，可以帮助系统提炼新闻重点并判断其潜在影响。本文系统通过新闻分析 Agent 对新闻标题、摘要和来源信息进行处理，并在无法获得真实新闻时明确提示“没有找到相关的新闻”或使用 Mock 回退提示。

\subsection{文献评述}

综上所述，已有研究说明大语言模型和多 Agent 协作机制在金融分析中具有较大潜力。TradingAgents 等框架展示了多角色协作在交易决策中的应用价值，但此类系统通常结构较复杂，完整复现和部署成本较高。与此同时，面向普通用户的轻量级股票辅助分析系统仍具有实践意义。基于此，本文结合 A 股和美股输入场景，设计并实现 QuantMind Agent 系统，通过九个核心 Agent 完成从数据获取、技术分析、新闻分析、基本面分析、舆情分析、多空研究、风险控制到最终决策展示的完整流程。

\section{研究设计}

\subsection{系统需求分析}

QuantMind Agent 的目标是实现一个可运行、可访问、可解释的股票辅助分析系统。根据项目目标，系统主要需求如下：

\begin{enumerate}
    \item 支持用户输入股票名称或股票代码，例如“贵州茅台”“600519”“Apple”“AAPL”等；
    \item 支持获取行情数据、新闻数据和基本面数据，并能够在接口不可用时进行回退；
    \item 支持技术分析、新闻分析、基本面分析、舆情分析、多空研究、风险控制和最终交易决策；
    \item 支持在网页中展示 Agent 执行进度和分析结果；
    \item 支持数据可信度提示和风险保护，避免基于占位数据输出高置信交易建议；
    \item 系统输出应具有可解释性，即不仅给出结果，还要给出依据和风险提示。
\end{enumerate}

\subsection{系统总体架构}

系统整体采用前后端结合的结构。后端使用 Python 和 FastAPI 实现，负责股票识别、数据获取、Agent 工作流执行和接口输出；前端使用 HTML、CSS 和 JavaScript 实现，负责用户输入、进度展示和结果卡片渲染。系统部署在 Zeabur 平台，用户可以直接通过浏览器访问。

系统总体流程如下：

\begin{center}
\begin{tabular}{c}
用户输入股票名称或代码 \\
$\downarrow$ \\
股票识别与输入校验 \\
$\downarrow$ \\
行情数据、新闻数据与基本面数据获取 \\
$\downarrow$ \\
技术分析 Agent $\rightarrow$ 新闻分析 Agent $\rightarrow$ 基本面分析 Agent $\rightarrow$ 舆情分析 Agent \\
$\downarrow$ \\
多头研究员 Agent $\rightarrow$ 空头研究员 Agent $\rightarrow$ 研究经理 Agent \\
$\downarrow$ \\
风险控制 Agent $\rightarrow$ 交易决策 Agent \\
$\downarrow$ \\
生成结构化分析结果 \\
$\downarrow$ \\
Web 页面展示最终决策和各 Agent 报告
\end{tabular}
\end{center}

\subsection{多 Agent 工作流程设计}

系统核心工作流位于 \texttt{quantmind/graph/workflow.py}。在工作流开始时，系统首先创建 AgentState 状态对象，并写入股票代码、分析日期、行情数据、新闻数据和基本面数据。之后系统依次调用九个 Agent：技术分析 Agent、新闻分析 Agent、基本面分析 Agent、舆情分析 Agent、多头研究员 Agent、空头研究员 Agent、研究经理 Agent、风险控制 Agent 和交易决策 Agent。

该流程可以表示为：

\begin{center}
\begin{tabular}{lll}
\toprule
步骤 & Agent & 主要任务 \\
\midrule
1 & 技术分析 Agent & 计算 MA5、MA10、最新价格、成交量变化，并判断技术趋势 \\
2 & 新闻分析 Agent & 根据新闻标题和摘要判断新闻情绪 \\
3 & 基本面分析 Agent & 根据 ROE、利润增长、估值、负债率等指标判断基本面强弱 \\
4 & 舆情分析 Agent & 判断市场情绪、关注热度和多空分歧程度 \\
5 & 多头研究员 Agent & 从已有报告中提取上涨或买入理由，形成多头观点 \\
6 & 空头研究员 Agent & 从已有报告中提取风险、下跌或观望理由，形成空头观点 \\
7 & 研究经理 Agent & 汇总多头与空头观点，形成研究层面的综合结论 \\
8 & 风险控制 Agent & 综合多维报告，输出风险等级、建议仓位和止损比例 \\
9 & 交易决策 Agent & 综合前面报告，输出 BUY、HOLD、WAIT 或 SELL \\
\bottomrule
\end{tabular}
\end{center}

与 TradingAgents 中复杂的多团队协作相比，本文系统采用固定串行流程，主要考虑到课程项目需要保证可运行性、可理解性和部署稳定性。当前系统没有实现完整的多轮辩论和基金经理审批流程，但已经通过多头研究员、空头研究员和研究经理三个 Agent 实现了轻量级观点对比与汇总，为后续扩展为更复杂的辩论式 Agent 架构奠定基础。

\subsection{数据获取与处理设计}

系统数据主要分为行情数据、新闻数据和基本面数据三类。行情数据可通过 AkShare 或 Tushare 获取，新闻数据可通过 AkShare 或 Alpha Vantage 获取，基本面数据由 \texttt{FundamentalDataProvider} 统一封装并根据股票市场类型选择可用数据源。为了保证系统在外部接口不可用时仍能运行，项目保留 Mock 数据或空指标回退方案。

在数据处理方面，系统会将行情数据整理为收盘价序列和成交量序列，供技术分析 Agent 计算指标。新闻数据则整理为标题、摘要、来源和链接等字段，供新闻分析 Agent 和舆情分析 Agent 分别判断新闻事件影响与市场情绪热度。基本面数据会整理为 ROE、利润增长率、营收增长率、利润率、PE、负债率、产权比率等指标，供基本面分析 Agent 进行财务质量判断。当前系统尤其重视数据可信度，如果数据源显示为 Mock 或 fallback Mock，最终交易决策 Agent 会触发保护逻辑，将最终动作降级为 WAIT，并将建议仓位设置为 0。

\subsection{核心 Agent 模块设计}

\subsubsection{技术分析 Agent}

技术分析 Agent 位于 \texttt{quantmind/agents/technical\_agent.py}。该模块首先检查行情数据长度是否足够，如果不足，则输出中性判断。若数据满足要求，则计算 MA5、MA10、最新价格和成交量变化率。其基本判断逻辑为：当最新价格高于 MA5 且 MA5 高于 MA10 时，认为技术面偏强；当最新价格低于 MA5 且 MA5 低于 MA10 时，认为技术面偏弱；其他情况下认为技术面中性。如果价格趋势偏强且成交量明显放大，系统会进一步提高技术评分。

\subsubsection{新闻分析 Agent}

新闻分析 Agent 位于 \texttt{quantmind/agents/news\_agent.py}。该模块根据新闻标题和摘要判断新闻情绪。规则模式下，系统会统计标题中正面关键词和负面关键词的出现情况。正面关键词包括“增长”“利好”“突破”“回购”“增持”“盈利”“创新高”等，负面关键词包括“下滑”“利空”“处罚”“减持”“亏损”“风险”“调查”等。如果正面词更多，则新闻情绪偏多；如果负面词更多，则新闻情绪偏空；若没有明显差异，则新闻情绪为中性。

\subsubsection{基本面分析 Agent}

基本面分析 Agent 位于 \texttt{quantmind/agents/fundamental\_agent.py}。该模块读取 \texttt{fundamental\_data} 中的 \texttt{metrics} 字段，并基于 ROE、利润同比增长率、营收同比增长率、利润率、PE、资产负债率和产权比率等指标判断公司基本面强弱。规则逻辑中，较高 ROE、正向利润增长、正向营收增长、较好利润率、估值未明显过高以及负债水平可控会增加偏多分数；利润增长为负、利润率偏低、PE 明显偏高或负债率较高则增加偏空分数。当可用财务指标不足时，系统会输出中性结论，并明确说明“没有可用的真实财务指标”或“可用指标较少”。如果配置了 DeepSeek API，基本面分析 Agent 还可以在规则基线基础上调用大语言模型生成中文摘要，但提示词要求模型只能基于已有财务字段分析，不得编造财报数据或行业事实。

\subsubsection{舆情分析 Agent}

舆情分析 Agent 位于 \texttt{quantmind/agents/sentiment\_agent.py}。该模块与新闻分析 Agent 的关注点不同：新闻分析 Agent 更强调新闻事件本身的利好或利空，而舆情分析 Agent 更关注市场情绪、关注热度和多空分歧程度。系统会根据新闻标题、摘要和来源统计正负面情绪词，输出 \texttt{sentiment}、\texttt{score}、\texttt{buzz\_score} 和 \texttt{disagreement\_score}。其中，\texttt{buzz\_score} 用于表示新闻数量带来的关注热度，\texttt{disagreement\_score} 用于表示正负情绪同时存在时的观点分歧程度。当新闻为空时，舆情分析 Agent 会输出中性、低热度和低分歧判断，避免在信息不足时给出过度明确的情绪结论。

\subsubsection{多头研究员 Agent}

多头研究员 Agent 位于 \texttt{quantmind/agents/research\_agent.py} 中的 \texttt{BullishResearchAgent}。该 Agent 的任务是站在多头研究员角度，从技术分析、新闻分析、基本面分析和舆情分析报告中寻找上涨、买入或正向支撑理由，并形成结构化多头观点。输出内容包括 \texttt{stance}、\texttt{confidence}、\texttt{thesis}、\texttt{key\_points} 和 \texttt{concerns}。如果已有报告中存在明确偏多信号，系统会将其整理为多头要点；如果存在偏空或中性信号，则会放入顾虑项中。该设计能够让系统在最终决策前先独立构造正向论证，而不是直接把所有信号简单相加。

\subsubsection{空头研究员 Agent}

空头研究员 Agent 同样位于 \texttt{quantmind/agents/research\_agent.py}，对应类为 \texttt{BearishResearchAgent}。该 Agent 的任务是站在空头或风险审查角度，从已有报告中寻找下跌、风险、观望或不确定性理由。与多头研究员 Agent 相对，空头研究员会把技术面偏弱、新闻面偏空、基本面不足、舆情负面或分歧较高等内容整理为空头要点，并把已有利多证据作为削弱空头观点的顾虑。该模块有助于避免系统只关注买入理由，使分析过程更接近真实投研中“正反两方论证”的思路。

\subsubsection{研究经理 Agent}

研究经理 Agent 位于 \texttt{quantmind/agents/research\_agent.py} 中的 \texttt{ResearchManagerAgent}。该模块负责汇总多头研究员和空头研究员的观点，并结合技术、新闻、基本面和舆情等基础分析报告，判断哪一方证据更有说服力。其输出包括 \texttt{conclusion}、\texttt{confidence}、\texttt{bullish\_summary}、\texttt{bearish\_summary}、\texttt{final\_summary} 和 \texttt{key\_evidence}。在规则逻辑中，研究经理会根据多头报告、空头报告和基础分析信号计算双方权重，若多头证据明显更强则输出偏多结论，若空头证据明显更强则输出偏空结论，若双方接近或信息不足则保持中性。研究经理结论随后会被风险控制 Agent 和交易决策 Agent 继续使用，成为最终辅助决策的重要中间依据。

\subsubsection{风险控制 Agent}

风险控制 Agent 位于 \texttt{quantmind/agents/risk\_agent.py}。该模块综合技术分析、新闻分析、基本面分析、舆情分析和研究经理结论计算风险评分。技术面、新闻面、基本面、舆情面和研究经理结论偏多会适当降低风险评分；相反，如果这些报告偏空，或者舆情分歧程度较高，则会提高风险评分。系统根据风险评分输出 low、medium、high 三类风险等级，并给出建议仓位和止损比例。该模块的作用是避免系统只关注买卖方向，而忽视仓位和风险约束。

\subsubsection{交易决策 Agent}

交易决策 Agent 位于 \texttt{quantmind/agents/decision\_agent.py}。该模块综合技术分析、新闻分析、基本面分析、舆情分析、研究经理结论和风险控制报告，输出最终辅助决策。系统支持四类动作：BUY 表示买入或试探性建仓，HOLD 表示已有仓位继续持有但不新增买入，WAIT 表示观望等待，SELL 表示卖出或规避风险。当前系统没有真实持仓上下文，因此非 BUY 动作的仓位通常设置为 0。

交易决策 Agent 还实现了重要的安全保护逻辑。当行情数据源包含 Mock 或 fallback 信息时，系统会将最终动作强制调整为 WAIT，置信度不超过 0.55，仓位为 0，并在摘要和风险提示中说明“未找到可用行情数据，无法基于真实行情给出买入或卖出判断”。

\subsection{Web 系统设计}

Web 后端位于 \texttt{web\_app.py}。系统提供三个主要接口：首页接口 \texttt{/}、股票校验接口 \texttt{/api/validate} 和流式分析接口 \texttt{/api/analyze/stream}。其中，\texttt{/api/analyze/stream} 使用 Server-Sent Events 技术向前端推送进度信息，使用户能够看到“已识别股票并获取基础数据”“技术分析完成”“新闻分析完成”“基本面分析完成”“舆情分析完成”“多头研究完成”“空头研究完成”“研究经理汇总完成”“风险控制完成”“最终决策生成”等过程。

前端页面位于 \texttt{web/static/index.html} 和 \texttt{web/static/app.js}。首页提供股票输入框和示例按钮；分析页通过进度条展示 Agent 执行状态，并在分析完成后展示最终决策卡片、技术分析卡片、新闻分析卡片、基本面分析卡片、舆情分析卡片、多空研究卡片、研究经理结论卡片和风险控制卡片。最终决策优先展示，包括动作、置信度、建议仓位、分析日期、风险提示和免责声明。

\subsection{风险控制与安全提示设计}

股票分析系统如果缺少风险提示，容易造成用户误解。因此，QuantMind Agent 在多个层面加入安全设计。首先，系统在最终结果中明确说明报告仅用于研究和学习，不构成投资建议。其次，风险控制 Agent 会综合技术、新闻、基本面、舆情和研究经理结论输出风险等级、建议仓位和止损比例。再次，当外部数据源不可用并回退到 Mock 数据时，系统会降低判断确信度，并强制最终决策为 WAIT。最后，LLM 调用失败时系统会回退到规则逻辑，保证系统不会因为模型接口异常而中断。

\section{实验分析}

\subsection{实验环境}

本文系统在本地开发环境和线上部署环境中进行测试。本地开发环境为 Windows 11，主要使用 Python 作为后端开发语言，FastAPI 作为 Web 服务框架，HTML、CSS 和 JavaScript 作为前端页面技术。数据接口包括 AkShare、Tushare 和 Alpha Vantage，大语言模型接口预留 DeepSeek。线上部署平台为 Zeabur，用户可以通过浏览器访问系统演示页面。

\begin{center}
\begin{tabular}{ll}
\toprule
项目 & 内容 \\
\midrule
后端语言 & Python \\
Web 框架 & FastAPI \\
前端技术 & HTML、CSS、JavaScript \\
行情数据源 & AkShare、Tushare、Mock 回退 \\
新闻数据源 & AkShare、Alpha Vantage、Mock 回退 \\
基本面数据源 & AkShare、Alpha Vantage、空指标回退 \\
LLM 接口 & DeepSeek，可失败回退到规则逻辑 \\
部署平台 & Zeabur \\
\bottomrule
\end{tabular}
\end{center}

\subsection{实验对象与实验流程}

为了验证系统功能，本文选取网页示例中常见的股票输入进行测试，包括“贵州茅台”“600519”“AAPL”和“特斯拉”等。系统支持股票名称和股票代码两种输入方式。若输入能够被识别，系统进入分析页面；若输入不支持，则前端提示当前内容不可查询。

实验流程如下：首先，用户在首页输入股票名称或代码；其次，系统调用 \texttt{/api/validate} 接口进行股票识别；再次，系统进入分析页面并通过 \texttt{/api/analyze/stream} 接口执行多 Agent 工作流；最后，前端接收分析结果并展示最终决策和各 Agent 报告。

\subsection{功能测试结果}

系统功能测试主要围绕输入识别、Agent 执行、结果展示和异常保护四个方面展开。测试结果如表所示。

\begin{center}
\begin{tabular}{lll}
\toprule
测试内容 & 预期结果 & 实际表现 \\
\midrule
股票输入识别 & 支持名称和代码输入 & 可识别 A 股和美股示例输入 \\
进度展示 & 显示 Agent 执行过程 & 前端进度条按步骤更新 \\
技术分析 & 输出技术趋势和指标 & 可展示 MA5、MA10、最新价和成交量变化 \\
新闻分析 & 输出新闻情绪和标题 & 可展示新闻摘要或无相关新闻提示 \\
基本面分析 & 输出基本面信号和财务指标 & 可展示 ROE、利润增长、估值、负债率等指标或数据不足提示 \\
舆情分析 & 输出市场情绪、热度和分歧 & 可展示 buzz\_score、disagreement\_score 和舆情来源 \\
多空研究 & 输出多头、空头和研究经理观点 & 可展示多头理由、空头风险和综合研究结论 \\
风险控制 & 输出风险等级和仓位 & 可展示风险等级、建议仓位和止损比例 \\
最终决策 & 输出 BUY/HOLD/WAIT/SELL & 可展示动作、置信度、仓位和风险提示 \\
异常保护 & 数据不可用时安全回退 & Mock 数据触发 WAIT 和 0 仓位保护 \\
\bottomrule
\end{tabular}
\end{center}

从测试结果可以看出，系统能够完成课程项目所需的主要功能。用户可以通过网页输入股票，系统能够按顺序执行九个 Agent，并在页面上展示结构化结果。

\subsection{案例分析}

以“贵州茅台”或“600519”为例，用户在首页输入股票名称或代码后，系统首先将输入识别为对应股票代码，然后进入分析页面。页面显示分析进度，依次完成基础数据准备、技术分析、新闻分析、基本面分析、舆情分析、多头研究、空头研究、研究经理汇总、风险控制和最终决策。

在技术分析阶段，系统根据最近价格序列计算 MA5、MA10、最新价格和成交量变化率。如果最新价格站上短期与中期均线，则系统倾向于给出偏多技术信号；如果跌破均线，则倾向于偏空；如果关系不明确，则保持中性。在新闻分析阶段，系统根据新闻标题和摘要判断新闻情绪，如果没有找到相关新闻，则明确说明“没有找到相关的新闻”，避免编造新闻依据。在基本面分析阶段，系统根据可用财务指标判断公司质量和估值压力；如果财务指标不足，则输出中性和数据不足提示。在舆情分析阶段，系统进一步判断市场关注热度和多空分歧程度。随后，多头研究员和空头研究员分别从正反两个角度整理证据，研究经理 Agent 对多空观点进行汇总。在风险控制阶段，系统综合多维报告给出风险等级、建议仓位和止损比例。在最终决策阶段，系统综合技术、新闻、基本面、舆情、研究经理结论和风险控制报告输出 BUY、HOLD、WAIT 或 SELL，并给出置信度、仓位和风险提示。

该案例体现了系统的两个特点。第一，分析过程是分步骤的，用户可以看到每个 Agent 的输出，而不是只看到一个最终结论。第二，系统结果具有可解释性，最终决策会同时展示技术依据、新闻依据、基本面依据、舆情依据、多空研究依据和风险控制说明。

\subsection{实验结果总结}

实验表明，QuantMind Agent 能够完成从用户输入到结果展示的完整流程，适合作为课程项目中的股票辅助分析系统。系统优势主要体现在以下几个方面：一是多 Agent 分工明确，系统结构清晰；二是新增基本面、舆情和多空研究模块后，分析维度更加完整；三是输出结果可解释，便于用户理解分析依据；四是 Web 页面操作简单，能够直观展示进度和结果；五是加入了 Mock 数据保护和风险提示，避免把测试数据误认为真实投资依据。

同时，系统也存在一定不足。首先，当前系统不是完整的自动交易系统，也没有接入真实账户交易。其次，技术指标较为基础，尚未加入 MACD、RSI、KDJ 等更多指标。再次，系统回测能力还不完善，不能直接证明策略长期收益。第四，当前多空研究仍是轻量级串行汇总，还没有实现复杂的多轮辩论、投票和审批机制。最后，新闻、行情和基本面数据质量受第三方接口影响，外部接口异常时需要依赖回退机制。

\section{结论与展望}

\subsection{研究结论}

本文设计并实现了 QuantMind Agent，一个基于多 Agent 的股票辅助分析系统。系统参考 TradingAgents 的多角色协作思想，但采用更轻量的固定串行流程，将股票分析过程划分为技术分析、新闻分析、基本面分析、舆情分析、多头研究、空头研究、研究经理汇总、风险控制和交易决策九个阶段。系统后端基于 Python 与 FastAPI 实现，前端基于 HTML、CSS 和 JavaScript 实现，并部署到 Zeabur 平台。用户可以通过网页输入股票名称或代码，系统自动执行分析流程并展示最终辅助决策结果。

通过实验测试可以看出，系统能够完成股票输入识别、Agent 流程执行、进度展示、结果卡片渲染和异常数据保护等功能。与单一规则判断相比，多 Agent 结构使分析过程更加清晰；与直接输出买卖建议相比，系统增加了风险控制和免责声明，使结果更适合作为学习和研究用途。

\subsection{研究不足}

本文系统仍存在一些不足。第一，系统当前主要用于辅助分析，不具备自动交易能力，也没有接入真实交易账户。第二，技术分析指标较少，主要包括 MA5、MA10、最新价格和成交量变化率，指标体系仍有扩展空间。第三，回测功能不完善，尚不能从长期历史数据角度验证决策效果。第四，多 Agent 协作方式仍以固定串行为主，多头研究、空头研究和研究经理模块属于轻量级观点汇总，还没有实现 TradingAgents 中更复杂的多轮辩论和审批机制。第五，新闻、行情和基本面数据依赖第三方接口，接口稳定性会影响系统分析质量。

\subsection{未来展望}

后续可以从以下方向继续改进。首先，可以增加 MACD、RSI、KDJ、布林带等更多技术指标，使技术分析更加全面。其次，可以完善回测模块，支持对历史数据进行策略效果验证。再次，可以进一步扩展多 Agent 辩论机制，在现有多头研究员、空头研究员和研究经理 Agent 基础上加入多轮交叉质询、投票和审批流程。第四，可以增强中文金融新闻情绪分析能力，提高 A 股场景下的新闻理解效果。第五，可以丰富基本面数据字段，加入现金流、行业估值分位和财报同比环比变化等指标。第六，可以加入投资组合管理功能，从单只股票分析扩展到多资产配置分析。最后，可以继续优化 Web 页面交互体验，使系统更适合展示和教学使用。

\newpage
\begin{thebibliography}{99}

\bibitem{tradingagents}
Xiao Y., Sun E., Luo D., Wang W. TradingAgents: Multi-Agents LLM Financial Trading Framework[EB/OL]. arXiv:2412.20138, 2025.

\bibitem{fingpt}
Yang H., Liu X. Y., Wang C. D. FinGPT: Open-Source Financial Large Language Models[EB/OL]. arXiv:2306.06031, 2023.

\bibitem{bloomberggpt}
Wu S., Irsoy O., Lu S., et al. BloombergGPT: A Large Language Model for Finance[EB/OL]. arXiv:2303.17564, 2023.

\bibitem{pixiu}
Xie Q., Han W., Zhang X., et al. PIXIU: A Large Language Model, Instruction Data and Evaluation Benchmark for Finance[EB/OL]. arXiv:2306.05443, 2023.

\bibitem{react}
Yao S., Zhao J., Yu D., et al. ReAct: Synergizing Reasoning and Acting in Language Models[EB/OL]. arXiv:2210.03629, 2023.

\bibitem{metagpt}
Hong S., Zhuge M., Chen J., et al. MetaGPT: Meta Programming for A Multi-Agent Collaborative Framework[EB/OL]. arXiv:2308.00352, 2024.

\bibitem{lopez}
Lopez-Lira A., Tang Y. Can ChatGPT Forecast Stock Price Movements? Return Predictability and Large Language Models[EB/OL]. arXiv:2304.07619, 2023.

\bibitem{finmem}
Yu Y., Li H., Chen Z., et al. FinMem: A Performance-Enhanced LLM Trading Agent with Layered Memory and Character Design[EB/OL]. arXiv:2311.13743, 2023.



\end{thebibliography}



\end{document}